% ============================================================================
% Unicode -> LaTeX fallback header for xelatex/lualatex/pdflatex.
%
% WHY: pandoc + xelatex with a text font (e.g. Liberation Serif) SILENTLY DROPS
% Unicode math/Greek glyphs the font lacks ("Missing character" warnings) — the
% PDF builds without error but symbols vanish (e.g. "D∩B" renders as "DB",
% "10⁻¹⁴" renders as "10¹⁴"). This maps the at-risk glyphs to LaTeX commands so
% the source .md can stay in clean Unicode (good for DOCX export and readability)
% while the PDF still renders every symbol.
%
% USAGE: pandoc paper.md ... --include-in-header=papers/_assets/unicode_fallback.tex
%   (wrapped by scripts/build_paper_pdf.py, which also diagnoses missing glyphs)
%
% MAINTENANCE: when scripts/build_paper_pdf.py reports a still-missing glyph,
% add one line:  \newunicodechar{<glyph>}{<latex>}  in the matching block below.
% Latin-1 glyphs that text fonts already provide (× ± ° ÷ – — −) are intentionally
% NOT remapped, to avoid changing their text-mode spacing/rendering.
% ============================================================================
\usepackage{amsmath}
\usepackage{amssymb}
\usepackage{textcomp}
\usepackage{newunicodechar}
% needspace : permet d'exiger N lignes libres avant un bloc (tables), pour
% interdire qu'un tableau soit coupé entre deux pages (papiers 1/2/3, 02/09).
\usepackage{needspace}
% array : les specs de colonnes calculées par scripts/table_layout.py utilisent
% >{\raggedright\arraybackslash}p{…} (colonnes à largeur exacte, non justifiées
% comme les minipages de pandoc). La préambule pandoc ne charge PAS array —
% sans lui, « Undefined control sequence \arraybackslash » au build (mesuré le
% 17/09 sur paper3). La sonde XeLaTeX de table_layout charge déjà array : les
% mesures et le rendu réel partagent ainsi exactement le même contexte.
\usepackage{array}
% booktabs : MÊME archetype de bug que array ci-dessus, mesuré le 2026-10-04
% sur le papier 5 — pandoc n'émet `\usepackage{longtable,booktabs}` que si son
% AST contient au moins une table ; or la politique de paper_fix_columns
% (2026-10-04) convertit TOUTES les tables pipe en blocs {=latex} — AST sans
% table, préambule sans booktabs, et le premier \toprule des longtables
% converties meurt en « Undefined control sequence » (xelatex ligne 723 du
% .tex intermédiaire, build EN du papier 5). fr_tables.tex chargeait déjà
% longtable en défense ; booktabs rejoint ici array, dans l'en-tête universel.
% Idempotent : sans effet quand pandoc a déjà chargé le paquet.
\usepackage{booktabs}

% --- Règle du 2026-09-19 (demande Guillaume) : AUCUN titre de section seul en bas de page.
% Valeurs MESURÉES après itérations (3 lignes : deux fuites — chaîne H2→H3 et needspace de
% tableau après titre ; 8/5/4 : « 4. Résultats » encore orphelin en p6 du papier 3 FR, la
% hauteur réelle d'un \section ≈ 3 lignes mange la marge de la chaîne 8→5). Marge prise :
% section 12 ≥ subsection 7 + hauteur section 3 + 2 ; subsection 7 ≥ subsubsection 5 + 2.
% Les \needspace de tableaux placés APRÈS un titre dans les .md ont été déplacés AVANT
% (scripts/paper1_move_needspace.py) — la garde TeX ne peut pas deviner ce qui suit.
% etoolbox : \pretocmd sans changer la mise en forme des titres.
\usepackage{etoolbox}
% \needspaceHeading fige la définition ORIGINALE de \needspace : fr_tables.tex (builds FR /
% --compact-tables) redéfinit \needspace pour poser un drapeau global nsExplicit devant les
% longtables — si la garde des titres passait par cette version redéfinie, elle ARMERAIT le
% drapeau et la longtable suivante SAUTERAIT sa garde anti-coupe automatique (régression).
\let\needspaceHeading\needspace
\pretocmd{\section}{\needspaceHeading{12\baselineskip}}{}{}
\pretocmd{\subsection}{\needspaceHeading{7\baselineskip}}{}{}
\pretocmd{\subsubsection}{\needspaceHeading{4\baselineskip}}{}{}

% --- set theory / logic --------------------------------------------------------
\newunicodechar{∩}{\ensuremath{\cap}}
\newunicodechar{∪}{\ensuremath{\cup}}
\newunicodechar{⊂}{\ensuremath{\subset}}
\newunicodechar{⊃}{\ensuremath{\supset}}
\newunicodechar{⊆}{\ensuremath{\subseteq}}
\newunicodechar{⊇}{\ensuremath{\supseteq}}
\newunicodechar{∈}{\ensuremath{\in}}
\newunicodechar{∉}{\ensuremath{\notin}}
\newunicodechar{∋}{\ensuremath{\ni}}
\newunicodechar{∅}{\ensuremath{\varnothing}}
\newunicodechar{∖}{\ensuremath{\setminus}}
\newunicodechar{∀}{\ensuremath{\forall}}
\newunicodechar{∃}{\ensuremath{\exists}}
\newunicodechar{∧}{\ensuremath{\wedge}}
\newunicodechar{∨}{\ensuremath{\vee}}
\newunicodechar{¬}{\ensuremath{\neg}}

% --- relations -----------------------------------------------------------------
\newunicodechar{≤}{\ensuremath{\leq}}
\newunicodechar{≥}{\ensuremath{\geq}}
\newunicodechar{≠}{\ensuremath{\neq}}
\newunicodechar{≈}{\ensuremath{\approx}}
\newunicodechar{≡}{\ensuremath{\equiv}}
\newunicodechar{∝}{\ensuremath{\propto}}
\newunicodechar{≪}{\ensuremath{\ll}}
\newunicodechar{≫}{\ensuremath{\gg}}
\newunicodechar{∼}{\ensuremath{\sim}}
\newunicodechar{≃}{\ensuremath{\simeq}}

% --- operators / analysis ------------------------------------------------------
\newunicodechar{⋅}{\ensuremath{\cdot}}
\newunicodechar{∘}{\ensuremath{\circ}}
\newunicodechar{∗}{\ensuremath{\ast}}
\newunicodechar{∞}{\ensuremath{\infty}}
\newunicodechar{∂}{\ensuremath{\partial}}
\newunicodechar{∇}{\ensuremath{\nabla}}
\newunicodechar{√}{\ensuremath{\surd}}
\newunicodechar{∑}{\ensuremath{\sum}}
\newunicodechar{∏}{\ensuremath{\prod}}
\newunicodechar{∫}{\ensuremath{\int}}

% --- arrows --------------------------------------------------------------------
\newunicodechar{→}{\ensuremath{\rightarrow}}
\newunicodechar{←}{\ensuremath{\leftarrow}}
\newunicodechar{↔}{\ensuremath{\leftrightarrow}}
\newunicodechar{⇒}{\ensuremath{\Rightarrow}}
\newunicodechar{⇐}{\ensuremath{\Leftarrow}}
\newunicodechar{⇔}{\ensuremath{\Leftrightarrow}}
\newunicodechar{↦}{\ensuremath{\mapsto}}

% --- Greek lowercase -----------------------------------------------------------
\newunicodechar{α}{\ensuremath{\alpha}}
\newunicodechar{β}{\ensuremath{\beta}}
\newunicodechar{γ}{\ensuremath{\gamma}}
\newunicodechar{δ}{\ensuremath{\delta}}
\newunicodechar{ε}{\ensuremath{\varepsilon}}
\newunicodechar{ζ}{\ensuremath{\zeta}}
\newunicodechar{η}{\ensuremath{\eta}}
\newunicodechar{θ}{\ensuremath{\theta}}
\newunicodechar{ι}{\ensuremath{\iota}}
\newunicodechar{κ}{\ensuremath{\kappa}}
\newunicodechar{λ}{\ensuremath{\lambda}}
\newunicodechar{μ}{\ensuremath{\mu}}
\newunicodechar{ν}{\ensuremath{\nu}}
\newunicodechar{ξ}{\ensuremath{\xi}}
\newunicodechar{π}{\ensuremath{\pi}}
\newunicodechar{ρ}{\ensuremath{\rho}}
\newunicodechar{σ}{\ensuremath{\sigma}}
\newunicodechar{τ}{\ensuremath{\tau}}
\newunicodechar{υ}{\ensuremath{\upsilon}}
\newunicodechar{φ}{\ensuremath{\varphi}}
\newunicodechar{χ}{\ensuremath{\chi}}
\newunicodechar{ψ}{\ensuremath{\psi}}
\newunicodechar{ω}{\ensuremath{\omega}}

% --- Greek uppercase -----------------------------------------------------------
\newunicodechar{Γ}{\ensuremath{\Gamma}}
\newunicodechar{Δ}{\ensuremath{\Delta}}
\newunicodechar{Θ}{\ensuremath{\Theta}}
\newunicodechar{Λ}{\ensuremath{\Lambda}}
\newunicodechar{Ξ}{\ensuremath{\Xi}}
\newunicodechar{Π}{\ensuremath{\Pi}}
\newunicodechar{Σ}{\ensuremath{\Sigma}}
\newunicodechar{Φ}{\ensuremath{\Phi}}
\newunicodechar{Ψ}{\ensuremath{\Psi}}
\newunicodechar{Ω}{\ensuremath{\Omega}}

% --- superscripts (the "10⁻⁴" case) --------------------------------------------
\newunicodechar{⁻}{\textsuperscript{\textminus}}
\newunicodechar{⁺}{\textsuperscript{+}}
\newunicodechar{⁰}{\textsuperscript{0}}
\newunicodechar{¹}{\textsuperscript{1}}
\newunicodechar{²}{\textsuperscript{2}}
\newunicodechar{³}{\textsuperscript{3}}
\newunicodechar{⁴}{\textsuperscript{4}}
\newunicodechar{⁵}{\textsuperscript{5}}
\newunicodechar{⁶}{\textsuperscript{6}}
\newunicodechar{⁷}{\textsuperscript{7}}
\newunicodechar{⁸}{\textsuperscript{8}}
\newunicodechar{⁹}{\textsuperscript{9}}

% --- subscripts ----------------------------------------------------------------
\newunicodechar{₀}{\textsubscript{0}}
\newunicodechar{₁}{\textsubscript{1}}
\newunicodechar{₂}{\textsubscript{2}}
\newunicodechar{₃}{\textsubscript{3}}
\newunicodechar{₄}{\textsubscript{4}}
\newunicodechar{₅}{\textsubscript{5}}
\newunicodechar{₆}{\textsubscript{6}}
\newunicodechar{₇}{\textsubscript{7}}
\newunicodechar{₈}{\textsubscript{8}}
\newunicodechar{₉}{\textsubscript{9}}
